\section{Limitations}
We identify a few limitations below:

\textbf{Vague labels for HumanEval samples}: We find that HumanEval specifications can often be vague with the inputs and outputs that are tested on. Therefore, some programs can be argued to be either correct or counterfeit. When manually inspecting programs and their scores, we find that base tests are too weak while EvalPlus tests are too strong. Therefore, for correctness, we use the criteria that the program must pass all base samples and at least 95\% of EvalPlus samples. However, this only affects a small fraction of samples and we do not believe changes any of our main claims (which are also supported by LeetCode and ODEX). 

\textbf{Filter for counterfeit samples}: In this work, we use a relatively liberal filter for counterfeit samples that consists of mostly basic syntax and/or correctness checks. While we believe our results would hold for slight alterations of our filter, we do not assess this.

\textbf{Nature of counterfeit samples}: The scope of this work is limited to counterfeit samples that are generated by sampling from a natural language description. It is unclear how these samples differ from human-written incorrect samples or samples constructed in a different way, for example by synthetically injecting bugs into correct samples as in HumanEvalFix \citep{muennighoff2023octopack}.

\textbf{Dataset and prompting variation}: While we make a best-effort attempt to use standardized prompts that lead to the best performance, evaluation has been found to be quite sensitive to the prompt and task format \citep{mizrahi2023state}. In addition, there is variation across the datasets generated by various models. We try to mitigate this by showcasing that our conclusions remain robust across a variety of datasets and models. 

\textbf{Other perspectives on code understanding}: Although the three tasks we evaluate capture important aspects of code understanding, our claims do not necessarily extrapolate to other aspects of code understanding such as code summarization, translation, or optimization. We believe that other dimensions of code understanding are equally important and encourage future evaluation beyond the tasks we present here.

\textbf{Limited results for GPT-3.5 and GPT-4}: All our counterfeit samples are generated from CodeLlama, DeepSeekInstruct, and StarCoder, so it is unknown whether the same insights apply to GPT-3.5 and GPT-4 counterfeits. In addition, due to budget constraints, we only evaluate these two models on a limited subset of our counterfeit datasets, decreasing the statistical significance of our results on these models. 